\section{Chapter 4}

\subsection{Buffer Size in Routers}

Traditional wisdom suggests:

\begin{equation*}
    B = RTT \cdot C
\end{equation*}

However, more recent research suggests:

\begin{equation*}
    B = \frac{RTT \cdot C}{\sqrt N}
\end{equation*}

where
\begin{itemize}
    \item $RTT$ is the round-trip time [sec],
    \item $C$ is the link capacity [bits/sec],
    \item $N$ is the \emph{independent} TCP flows sharing the link.
\end{itemize}

Based on more recent research, this formula applies when $N$ is large such as in the network core.

\subsection{Weighted Fair Queuing (WFQ)}

In output port packet scheduling, each class $i$ is guaranteed to receive a fraction of the service, given by

\begin{equation*}
    \frac{w_i}{\sum_{j=1}^N w_j}
\end{equation*}

where
\begin{itemize}
    \item $w_i$ is the weight of class $i$,
    \item $w_j$ is the weight of class $j$,
    \item $N$ is the number of classes with at least one packet in the queue.
\end{itemize}